% !TeX root = ../thuthesis-example.tex

% 中英文摘要和关键字

\begin{abstract}
中性粒细胞在执行先天免疫效应过程中长期面临强氧化应激、活跃膜运输与膜重塑以及颗粒/溶酶体高负荷等多重压力，其关键调控分子往往呈现明显的情境依赖性。受限于原代中性粒细胞寿命短、遗传操作困难，传统高通量筛选难以系统识别此类依赖蛋白。针对这一问题，本研究基于蛋白语言模型构建了 Human-level 与 Immune-level 双层必需性评估体系，用于区分基础生存依赖与免疫情境条件性依赖。本文所研究的免疫情境，在生物学解释与实验验证层面具体落脚于急性炎症应答中中性粒细胞的氧化爆发-颗粒/溶酶体高负荷活化状态。研究采用 ProtT5 提取序列高维表征，并构建基于单头注意力的判别架构，在独立测试集上取得 AUROC 0.9272 和 AUPRC 0.8245。基于双层 PES 评分及 $\Delta PES$ 比较，本研究筛得一批免疫情境候选必需蛋白；这些蛋白整体呈现正电荷偏移、分子量轻量化和疏水核心强化等特征，与中性粒细胞在酸性吞噬体、高氧化和高膜周转条件下的功能需求相一致。以 LyLAP (PLBD1) 为代表的注意力图谱与机制分析进一步支持了模型对 Immune-level 依赖信号的识别能力。在此基础上，ATP6V1B2 被识别为高置信候选靶点，其 Immune-level PES（0.952）显著高于 Human-level PES（0.007），提示其可能参与氧化爆发-颗粒/溶酶体高负荷活化状态下的酸化动力学调控。实验上，本文选用 HL-60 细胞作为中性粒细胞的体外替代体系，并通过 PMA 刺激模拟上述状态，从而近似再现氧化爆发增强及溶酶体/酸化负荷升高等关键特征。综合来看，本研究建立了从计算预测到机制解释的筛选框架，为发现具有细胞选择性的免疫调节靶点提供了新的技术路径。

  % 关键词用“英文逗号”分隔，输出时会自动处理为正确的分隔符
  \thusetup{
    keywords = {中性粒细胞, ATP6V1B2蛋白, 深度学习, 靶点筛选, 情境必需性},
  }
\end{abstract}

\begin{abstract*}
Neutrophils function under persistent oxidative stress, active membrane trafficking and remodeling, and heavy granule/lysosomal burden during innate immune responses, making their key regulators strongly context dependent. Because primary neutrophils are short-lived and difficult to manipulate genetically, conventional high-throughput screening cannot systematically identify these dependency proteins. To address this problem, we established a dual-layer essentiality evaluation framework based on protein language models to distinguish basal survival dependency (Human-level) from immune-context conditional dependency (Immune-level). In this study, the biological interpretation and experimental approximation of this immune context are focused on an activated neutrophil state during acute inflammation, characterized by enhanced oxidative burst and elevated granule/lysosomal load.

Using ProtT5 embeddings and a single-head attention classifier, the model achieved an AUROC of 0.9272 and an AUPRC of 0.8245 on an independent test set. Comparative analysis of Protein Essentiality Scores (PES) and $\Delta$PES identified a set of candidate essential proteins associated with the immune context. These proteins showed increased basicity, reduced molecular size, and reinforced hydrophobic cores, consistent with functional constraints imposed by acidic phagosomes, oxidative stress, and high membrane turnover in neutrophils. Attention-based and mechanistic analysis of LyLAP (PLBD1) further supported the model's ability to capture Immune-level dependency signals.

ATP6V1B2 emerged as a high-confidence target, with an Immune-level PES of 0.952 versus a Human-level PES of 0.007, suggesting a potential role in acidification dynamics in the oxidative-burst, granule/lysosomal high-burden activated state. For experimental approximation, HL-60 cells were used as an in vitro surrogate for neutrophils, and PMA stimulation was applied to reproduce key features of this state. Overall, this study establishes a framework linking computational prediction with mechanistic interpretation and provides a practical strategy for identifying selective immunomodulatory targets.

  % Use comma as separator when inputting
  \thusetup{
    keywords* = {Neutrophils, ATP6V1B2, Deep Learning, Target Screening, Context-specific Essentiality},
  }
\end{abstract*}
